\documentclass[11pt]{article}
\usepackage[margin=1in]{geometry}
\usepackage[OT1]{fontenc}
\usepackage{amsmath,amssymb,amsthm}
\usepackage{microtype}
\usepackage[unicode,colorlinks=true,linkcolor=blue,citecolor=blue,urlcolor=blue]{hyperref}
\renewcommand{\rmdefault}{cmr}
\renewcommand{\sfdefault}{cmss}
\renewcommand{\ttdefault}{cmtt}
\setlength{\parindent}{0pt}
\setlength{\parskip}{5pt plus 1pt}
\setlength{\emergencystretch}{3em}
\allowdisplaybreaks[1]

\newtheorem{theorem}{Theorem}
\newtheorem{lemma}[theorem]{Lemma}
\newtheorem{proposition}[theorem]{Proposition}
\theoremstyle{definition}
\newtheorem{definition}[theorem]{Definition}
\newtheorem{example}[theorem]{Example}
\renewcommand{\qedsymbol}{\ensuremath{\square}}
\newcommand{\N}{\mathbb N}
\newcommand{\PA}{\mathsf{PA}}
\newcommand{\Tr}{\mathsf{Tr}}
\newcommand{\FV}{\operatorname{FV}}
\newcommand{\ucl}{\operatorname{ucl}}
\newcommand{\code}[1]{\ulcorner #1\urcorner}
\newcommand{\num}[1]{\overline{#1}}
\newcommand{\seqcode}[1]{\operatorname{code}(#1)}

\title{An explicit ill-founded hierarchy\protect\\of self-referential theories}
\author{}
\date{27 October 2026}
\hypersetup{pdftitle={An explicit ill-founded hierarchy of self-referential theories}}

\begin{document}
\maketitle

\begin{abstract}
Alexander proved that some ill-founded recursively enumerable partial order
supports a true family of self-referential theories in which every member
expresses a particular code of every member and the truth of its subordinates.
His proof establishes existence without specifying a witnessing order. We give
a recursive order and an effective construction of such a family. The order is
the reverse-extension relation on a tree of finite arithmetic truth tables,
existential witnesses and diagonal values. The tree has infinite branches,
but none computable from arithmetic truth. For Alexander's fixed-point family,
unsound members can be enumerated using arithmetic truth; any unsound member
would therefore generate a forbidden descending path. This replaces the
well-founded induction in his constructive theorem. The same order works
uniformly for every closed-r.e.-generic background, including arithmetic for
the full modal language. We explain its compatibility with the obstruction to
computable successor hierarchies and its relation to the minimal introspection
bundle.
\end{abstract}

\section{The question and the construction}\label{sec:introduction}

A theory can describe a program that enumerates another theory's conclusions.
It can also assert that those conclusions are true. These are different kinds
of information: a program tells us what the other theory says, while a
truthfulness assertion licenses accepting what it says.

Alexander studies families of recursively enumerable theories that can express
both kinds of information about one another. In \emph{Self-referential
theories}~\cite{alexander}, Theorem~24 constructs true families on every
recursively enumerable well-founded partial order. Every member expresses a
particular code of every member, and it expresses the truth of the members
below it. Well-founded induction proves the constructed family true.

Theorem~79(1) of the same paper establishes that some ill-founded r.e. orders
also support this construction. Its proof compares the complexity of
well-foundedness with the expressive power of arithmetic truth. It does not
identify an order on which the construction succeeds. The present paper gives
one by a direct finite membership test.

We use \emph{recursive} and \emph{computable} interchangeably. An r.e. set is
a set that an ordinary program can enumerate, also called a computably
enumerable set. A family $(T_i)_{i\in\N}$ is uniformly r.e. when a single
program enumerates all pairs $(i,\sigma)$ with $\sigma\in T_i$. A strict
partial order is an irreflexive, transitive relation. We write $j\prec i$ when
$j$ is subordinate to $i$, so a descending path has the form
\[
 i_0\succ i_1\succ i_2\succ\cdots.
\]

Our language is ordinary arithmetic with effectively named operators
$K_0,K_1,\ldots$. The canonical interpretation reads $K_j\phi$ as an assertion
that $T_j$ entails the corresponding sentence. Section~\ref{sec:framework}
gives the precise convention. A family is \emph{true} if every one of its
axioms is true in this common standard interpretation.

\begin{theorem}[Explicit witnessing hierarchy]\label{thm:main}
There is an algorithmically specified recursive strict partial order
$\prec_*$ on $\N$ with the following properties.
\begin{enumerate}
\item The order is ill-founded, but it has no infinite descending path
computable from the set of true arithmetic sentences.
\item An effective fixed-point construction gives a true, uniformly r.e.
family $(T_i)_{i\in\N}$ in which every $T_i$ contains the Peano axioms for the
full modal language, a particular-code schema for every $T_j$, and the full
truth schema for every $T_j$ with $j\prec_*i$.
\item The same order works for every closed-r.e.-generic background in
Alexander's sense. From an index of that background, a total computable
compiler returns a successful fixed-point index.
\end{enumerate}
\end{theorem}

Here ``a particular code'' means a schema containing an actual numeral index,
rather than an assertion that some index exists. In Alexander's common-index
presentation, every $T_i$ contains, for the same natural number $n$,
\begin{equation}\label{eq:code-schema}
 \forall x\bigl(K_j\phi\ \leftrightarrow\
 \langle\num{\code{\phi}},\num j,x\rangle\in W_{\num n}\bigr),
 \qquad \FV(\phi)\subseteq\{x\}.
\end{equation}
The finite tuple notation in this arithmetic formula denotes a fixed
computable coding. The truth schema for a subordinate is
\begin{equation}\label{eq:truth-schema}
 \ucl(K_j\phi\to\phi),\qquad j\prec_*i,
\end{equation}
where $\ucl$ means universal closure. The numeral $\num n$ in
\eqref{eq:code-schema} is produced effectively; it is not selected by an
oracle that tests truth.

The construction separates two questions about a hierarchy. Its relation can
be computable even though no program can follow an infinite descending path
through it. Ill-foundedness asserts that some such path exists. Computable
ill-foundedness would additionally provide an algorithm for choosing its
successive members. Theorem~\ref{thm:main} uses the first possibility while
excluding the second, and even excludes descent computable from arithmetic
truth.

There are two main steps. First we show that the absence of a path computable
from arithmetic truth is sufficient for Alexander's fixed-point family to be
true. Second we specify a recursive tree with exactly the required path
obstruction. The subordination order is reverse proper extension of its
finite nodes. Sections~\ref{sec:order}--\ref{sec:assembly} assemble the two
steps, and Section~\ref{sec:comparison} explains the relation to the
successor-chain problem and the introspection result.

The argument below is a hand proof. It does not claim independent
mathematical review, formal verification or a determination of literature
priority. It establishes an explicit witness for the construction in
Theorem~79(1); it does not assert the corresponding result for the different
stratified construction in part~(2).

\section{Canonical theories and effective code knowledge}\label{sec:framework}

\subsection{The consequence relation and the standard interpretation}

Write $L=L_{\mathrm{PA}}(\omega)$ for Alexander's arithmetic language with operators
$K_i$. We use his base logic, in which interpretations of the operators
respect free variables, alphabetic variants and substitution of variables
\cite[Definition~2]{alexander}. Write $U\vdash\sigma$ for consequence in
that logic. The published logical input is its completeness and compactness
theorem: consequences of an r.e. theory can be enumerated effectively,
uniformly from an enumeration index of its axioms
\cite[Theorem~3]{alexander}. We use this actual consequence relation
throughout.

For an assignment $\eta$ of natural numbers to the free variables of $\phi$,
let $\phi^\eta$ be the sentence obtained by substituting the corresponding
numerals, renaming bound variables where necessary. For a family $U=(U_i)$,
its canonical structure $M_U$ has standard arithmetic as its first-order
part and satisfies
\begin{equation}\label{eq:canonical}
 M_U\models K_i\phi[\eta]
 \quad\Longleftrightarrow\quad U_i\vdash\phi^\eta.
\end{equation}
Alexander's substitution lemma verifies that
\begin{equation}\label{eq:substitution}
 M_U\models\phi[\eta]
 \quad\Longleftrightarrow\quad M_U\models\phi^\eta
\end{equation}
for every formula~\cite[Lemma~8]{alexander}. A true family means
$M_U\models U_i$ for every $i$. In that event every consequence of every
$U_i$ is also true in the same structure.

\subsection{Closure and admissible background axioms}

\begin{definition}[Closure and generic background]\label{def:generic}
A family $U$ is \emph{closed} if $\sigma\in U_i$ implies
$K_i\sigma\in U_i$ for all $i$. A uniformly r.e. background family $B$ is
\emph{closed-r.e.-generic} if $M_U\models B$ for every closed uniformly r.e.
family $U\supseteq B$.
\end{definition}

Closure here refers to the addition of knowledge assertions to the axiom
sets. It is distinct from closing an axiom set under all logical
consequences. The latter are already accounted for by $\vdash$.

Genericity gives a convenient way to include ordinary reasoning principles
without settling the desired hierarchy in advance. Once a closed r.e.
family extending $B$ has been constructed, its canonical interpretation
automatically satisfies $B$. The uniform family consisting of the Peano
axioms for $L$, including induction for its modal formulas, is one example
\cite[Lemmas~10 and~20]{alexander}. Section~\ref{subsec:introspection}
discusses a further example using the smaller introspection bundle.

\subsection{The effective fixed point}\label{subsec:fixed-point}

Fix an r.e. relation $\prec$ and an enumeration index $b$ of a background
family $B$. At this stage the relation need not be well-founded. For each
number $n$, let $G_i(n)$ consist of the background $B_i$, all instances of
\eqref{eq:code-schema}, and all instances of \eqref{eq:truth-schema} for
$j\prec i$. Put
\begin{equation}\label{eq:finite-closure}
 T_i(n)=\{K_i^k\sigma:\sigma\in G_i(n),\ k\in\N\},
 \qquad K_i^0\sigma=\sigma.
\end{equation}
This is precisely the smallest closed family containing those generators.
The definition is effective. Enumerate the background, dovetail the code
schemas over $j$ and $\phi$, enumerate subordinate truth schemas as pairs
appear in $\prec$, and enumerate all finite iterates $K_i^k$ of each
generator. No search for a true theory occurs in this enumeration.

Uniform consequence enumeration now gives a total computable function
$f_b$ on indices such that
\begin{equation}\label{eq:index-map}
 W_{f_b(n)}=
 \bigl\{\langle\code{\phi},j,m\rangle:
 \FV(\phi)\subseteq\{x\},\quad
 T_j(n)\vdash\phi(x\mid\num m)\bigr\}.
\end{equation}
Totality here means that the compiler always returns a program index. The
program it returns enumerates a set and can run forever, as an enumerator
normally does.

The effective recursion theorem supplies an index $n$ satisfying
\begin{equation}\label{eq:fixed-point}
 W_n=W_{f_b(n)}.
\end{equation}
Fix this $n$ and abbreviate $T=T(n)$ and $M=M_T$. Equation
\eqref{eq:fixed-point}, together with \eqref{eq:canonical}, makes every
code schema \eqref{eq:code-schema} true in $M$. This conclusion is
independent of well-foundedness. The remaining issue is the truth of the
subordinate truth schemas.

For later use, it is helpful to make the effectiveness of the index choice
explicit. Fix a numbering $(\varphi_x)$ of ordinary partial computable
functions. Let $\delta(x)$ be the effectively compiled index of the
enumerator that first computes $\varphi_x(x)$ and, if this returns $r$,
enumerates $W_r$. Parameter compilation makes $\delta$ a total computable
function. Compile an ordinary program $q$ for the total function
\[
 x\longmapsto f_b(\delta(x)),
\]
and output $n=\delta(q)$. Then
\begin{equation}\label{eq:explicit-compiler}
 W_n=W_{\varphi_q(q)}
     =W_{f_b(\delta(q))}=W_{f_b(n)}.
\end{equation}
All these operations are ordinary index compilation. The numerical output
depends on the chosen acceptable program numbering, just as any program's
index does. The procedure for obtaining it is fixed and requires no
uncomputable choice.

\section{How an unsound family would yield a path}\label{sec:soundness}

Let $A$ be the set of true arithmetic sentences in any fixed effective
coding. An $A$-computable procedure is an ordinary oracle computation that
may ask whether a specified arithmetic sentence is true. This oracle is an
external device in the proof. It is not added to the machines' language or
used by their enumerators.

The next lemma identifies why this particular oracle suffices.

\begin{lemma}[Canonical truth reduces to arithmetic truth]\label{lem:translation}
From an enumeration index of any uniformly r.e. family $U$, and from an
$L$-sentence $\sigma$, one can effectively produce an arithmetic sentence
$F_U(\sigma)$ such that
\begin{equation}\label{eq:arithmetic-translation}
 M_U\models\sigma
 \quad\Longleftrightarrow\quad\N\models F_U(\sigma).
\end{equation}
In particular, $A$ decides the truth of each individual canonical sentence,
uniformly in $U$ and $\sigma$.
\end{lemma}

\begin{proof}
Uniform consequence enumeration gives an arithmetic predicate
$\operatorname{Pr}_U(j,z)$ saying that $U_j$ entails the sentence with code
$z$. For example, it can assert the existence of a finite computation
witness in that enumeration.

Translate arithmetic atoms identically, and commute the translation with
Boolean operations and quantifiers. For a modal formula $K_j\psi$ with
free variables among $\vec x$, use the arithmetic formula saying that
$U_j$ entails the result of substituting the $\vec x$ numerals into
$\psi$. More explicitly, if
$\operatorname{Sub}_\psi(\vec x,z)$ is an arithmetic representation of
this computable substitution operation, set
\[
 F_U(K_j\psi)(\vec x)=
 \exists z\bigl(\operatorname{Sub}_\psi(\vec x,z)
                 \wedge\operatorname{Pr}_U(j,z)\bigr).
\]
This is an arithmetic formula. Induction on the outer construction of a
formula, using \eqref{eq:canonical} in the modal case, proves the claimed
equivalence. This is the translation used in Alexander's proof of
Theorem~79~\cite{alexander}.
\end{proof}

The modal case does not ask a truth oracle to evaluate the formula inside
the operator. It refers to a proof of that formula as a syntactic object.
For example, $K_iK_j(0=1)$ becomes an arithmetic assertion that $U_i$
entails the literal modal sentence $K_j(0=1)$. This is why arbitrary modal
nesting still gives an arithmetic translation.

\begin{lemma}[Downward propagation of unsoundness]\label{lem:local-step}
For the fixed-point family of Section~\ref{subsec:fixed-point}, assume $B$
is closed-r.e.-generic. Then
\begin{equation}\label{eq:local-step}
 \bigl(\forall j\prec i\;M\models T_j\bigr)
 \quad\Longrightarrow\quad M\models T_i.
\end{equation}
\end{lemma}

\begin{proof}
The background generators are true because $T$ is a closed r.e. extension
of $B$. The code generators are true by the fixed-point equation.

Suppose all subordinate theories are true. If
$M\models K_j\phi[\eta]$ for a subordinate $j$, then
$T_j\vdash\phi^\eta$. Since $M\models T_j$, it follows that
$M\models\phi^\eta$, and hence $M\models\phi[\eta]$ by the
substitution property. This verifies every subordinate truth generator.

Finally, every closure-generated axiom $K_i\rho$ is true in $M$: the
sentence $\rho$ belongs to $T_i$, so $T_i\vdash\rho$. This argument does
not assume that $\rho$ is true. It deals with every positive iterate in
\eqref{eq:finite-closure}. Thus all axioms of $T_i$ are true.
\end{proof}

This is the local induction step in the published well-founded argument.
We will use its contrapositive: if a member is unsound, at least one of its
subordinates is unsound. To turn that observation into a contradiction on
an ill-founded order, we need an effective way to find such a subordinate.

\begin{proposition}[Soundness criterion]\label{prop:criterion}
Let $\prec$ be an r.e. strict partial order with no $A$-computable infinite
descending path. For every closed-r.e.-generic background $B$, the
effectively compiled fixed-point family $T$ is true.
\end{proposition}

\begin{proof}
Let
\begin{equation}\label{eq:bad-set}
 D=\{i:(\exists\sigma\in T_i)\ M\not\models\sigma\}.
\end{equation}
This set is c.e. relative to $A$. Enumerate all axioms of all $T_i$, apply
Lemma~\ref{lem:translation} to each one, and use $A$ to check its truth.
Whenever a false axiom of $T_i$ appears, enumerate $i$ into $D$.

By Lemma~\ref{lem:local-step}, every $i\in D$ has a $j\in D$ with
$j\prec i$. Suppose $D$ is nonempty, and fix one $i_0\in D$. Given
$i_k$, dovetail the $A$-enumeration of $D$ with the enumeration of the
relation until a pair $(j,i_k)$ has appeared and $j$ has appeared in $D$.
Choose the first such discovery and put $i_{k+1}=j$. The local step
guarantees that this finite search terminates at every stage. Iterating it
computes an infinite descending path from $A$, contrary to the hypothesis.
Therefore $D$ is empty, so $T$ is true.
\end{proof}

Two effectiveness details matter here. First, the procedure only enumerates
unsound indices. It does not decide whether an entire r.e. theory is sound;
that would ask a stronger question than checking each of its individual
axioms. Consequently no jump of the truth oracle is being used in this
search. Second, the starting index $i_0$ can be a fixed numerical parameter
in the oracle program. The contradiction needs the existence of such a
program if $D$ is nonempty; it does not need a uniform method of finding
the first unsound theory in an arbitrary family.

The mechanism can be summarized as
\[
 \begin{aligned}
 &\text{one unsound member}\\
 &\quad\Longrightarrow\text{an $A$-enumerable set of unsound members}\\
 &\quad\Longrightarrow\text{an $A$-computable descending path}.
 \end{aligned}
\]
We next construct a recursive order where the last line is impossible,
although infinite descending paths exist.

\section{A recursive tree with inaccessible infinite branches}\label{sec:tree}

The tree will be built from finite records that propose answers to arithmetic
questions. Every infinite branch must eventually answer all those questions
correctly. A further coordinate makes it impossible for arithmetic truth
itself to compute an entire branch.

\subsection{The three coordinates}

Fix a recursive list $\theta_0,\theta_1,\ldots$, without repetitions, of all
closed arithmetic formulas written using atomic formulas, negation,
conjunction and existential quantification. One specific procedure is to list
finite strings in length and lexicographic order and retain the closed
well-formed formulas. The index of any given closed formula is computable by
searching this list. Other connectives are effective abbreviations.

For this section let $A(m)=1$ exactly when $\theta_m$ is true. This coding of
arithmetic truth is computably equivalent to the coding in the preceding
section. In particular, changing between them does not change which functions
are computable from arithmetic truth.

Fix a standard effective enumeration $(\Phi_e^X)$ of oracle programs, with
natural number inputs and outputs, together with a step-by-step universal
simulation. A candidate node of length $N$ is
\begin{equation}\label{eq:record}
 s=\bigl((b_0,w_0,d_0),\ldots,(b_{N-1},w_{N-1},d_{N-1})\bigr),
 \qquad b_m\in\{0,1\},\quad w_m,d_m\in\N.
\end{equation}
The entries have the following roles.

\begin{center}
\begin{tabular}{p{0.12\textwidth}p{0.35\textwidth}p{0.42\textwidth}}
\hline
Entry & Proposed information & What the finite checks enforce\\
\hline
$b_m$ & The truth value of $\theta_m$. & Arithmetic truth conditions whenever
the required entries are available.\\[3pt]
$w_m$ & A witness if $\theta_m$ is declared to be a true existential. & The
selected instance must be declared true once its index is available.\\[3pt]
$d_m$ & A value chosen against oracle program $m$. & It differs from the
output of $\Phi_m^b(m)$ whenever its halting computation is available.\\
\hline
\end{tabular}
\end{center}

The witness entry matters because a true existential sentence requires some
particular natural number witnessing it. Putting that number into the node
makes verification finite once the corresponding truth-table entry becomes
available. The diagonal entry serves a different purpose: after the entire
truth table has become correct, it prevents that truth table from computing
the rest of the branch.

\subsection{The complete finite membership test}\label{subsec:membership}

Reject a decoded record if any $b_m$ lies outside $\{0,1\}$. Otherwise apply
the following checks for every $m<N$.

\begin{enumerate}
\item If $\theta_m$ is atomic, require that $b_m$ is its actual truth value
in $\N$. Truth of a closed atomic arithmetic formula is directly computable.
\item If $\theta_m=\neg\theta_r$ and $r<N$, require $b_m=1-b_r$.
\item If $\theta_m=\theta_r\wedge\theta_t$ and $r,t<N$, require
$b_m=b_rb_t$.
\item If $\theta_m=\exists v\,\psi(v)$ and $b_m=1$, compute the index $r$
of the closed formula $\psi(\num{w_m})$. If $r<N$, require $b_r=1$.
\item If $\theta_m=\exists v\,\psi(v)$ and $b_m=0$, then for every $k<N$
compute the index $r$ of $\psi(\num k)$. Whenever $r<N$, require $b_r=0$.
\end{enumerate}

Substitution is capture-free, with the chosen bound variable replaced by a
numeral. The body of a closed existential has no other free variables, so
every substituted formula used in these tests is closed. Consequently each
inverse-index search terminates.

There is one further group of tests. For every $e<N$, simulate
$\Phi_e^b(e)$ for at most $N$ steps, using $b_0,\ldots,b_{N-1}$ to answer
its oracle queries. If a query has number at least $N$, defer that simulation:
there is no available answer. If the program halts within the allowed steps,
all its queries having been answered from the finite table, and its output
is $y$, require
\begin{equation}\label{eq:finite-diagonal}
 d_e\ne y.
\end{equation}

Let $P$ be the set of accepted finite records. Its membership test uses only
finite syntax calculations, evaluation of closed atomic formulas and bounded
program simulations. It never asks whether an arbitrary arithmetic sentence
is actually true. The proposed answers are input data to be checked for
finite consistency.

\begin{example}[An existential check]\label{ex:existential}
Let $a$ be the index of $0=0$, and let $c$ be the index of
$\exists v\,(v=0)$. These need not be the first entries of the fixed list.
Once a prefix contains both positions, the atomic test forces $b_a=1$.

If the prefix declares $b_c=1$ and chooses $w_c=0$, its positive existential
test agrees with $b_a=1$. If it instead chooses $w_c=1$, then at a sufficiently
long prefix it must declare $1=0$ true, contradicting an atomic test. If it
declares $b_c=0$, the negative test with $k=0$ eventually requires $b_a=0$,
which also contradicts the atomic test.

Thus the tests rule out both an incorrect positive witness and a false
negative answer. This second check is essential; merely checking witnesses
for declared true existentials would leave false negative answers unchecked.
\end{example}

\begin{example}[An unavailable oracle answer]
Suppose a simulated program asks for $b_q$ with $q\geq N$. Answering zero
would manufacture a computation that need not be the computation of any
eventual branch. The test therefore waits for a longer prefix. By contrast,
if the program has already halted using only entries below $N$, every
extension repeats that same finite computation. Its inequality
\eqref{eq:finite-diagonal} can be enforced permanently.
\end{example}

\subsection{Why the accepted records form a tree}

\begin{lemma}\label{lem:recursive-tree}
The set $P$ is recursive and closed under taking prefixes.
\end{lemma}

\begin{proof}
The membership algorithm has finitely many checks. All syntax operations
and inverse-index searches terminate, and each program simulation has a
specified finite step bound. Hence membership is decidable.

If a prefix fails a truth-table check, every extension fails that same
check: the relevant entries and the selected witness have not changed.
The negative existential checks include every smaller witness already
tested. If a prefix fails a diagonal check, the offending computation has
already halted using available oracle answers; every extension supplies
the same answers to those queries and has a larger simulation budget.
Thus failure at a prefix persists in every extension. Equivalently, every
prefix of an accepted record is accepted.
\end{proof}

The tree is a tree of finite sequences of natural number triples. There is
no finite or computable bound imposed on the witness and diagonal
coordinates. We will prove that it has a branch by specifying one
mathematically, rather than inferring that conclusion solely from the
existence of arbitrarily long finite nodes.

\subsection{Existence of a branch and the diagonal obstruction}

\begin{lemma}\label{lem:branches}
The tree $P$ has an infinite branch, but none of its infinite branches is
computable from $A$.
\end{lemma}

\begin{proof}
For existence, take the true arithmetic table $b_m=A(m)$. If $\theta_m$ is
a true existential sentence, let $w_m$ be its least true witness; otherwise
put $w_m=0$. Define
\begin{equation}\label{eq:branch-diagonal}
 d_e=\begin{cases}
 \Phi_e^A(e)+1,&\text{if }\Phi_e^A(e)\text{ halts},\\
 0,&\text{if it does not halt}.
 \end{cases}
\end{equation}
Every finite prefix of $(b,w,d)$ passes the truth-table tests. Whenever a
finite simulation halts using that prefix, it is the actual computation
with oracle $A$, so its output differs from $d_e$ by
\eqref{eq:branch-diagonal}. Thus every prefix belongs to $P$.

Now take any infinite branch $(b,w,d)$. We first show $b_m=A(m)$ for every
$m$. Induct on the number of logical connectives and quantifiers in the
closed formula $\theta_m$. Atomic truth is enforced directly. Negation
and conjunction have their usual truth values once the finitely many
required indices appear in a sufficiently long prefix.

For an existential $\theta_m=\exists v\,\psi(v)$, if $b_m=1$ the
selected witness $w_m$ has a substituted instance eventually checked to
have value $1$. By induction that instance is true, so the existential is
true. If $b_m=0$, each individual $k$ and the index of
$\psi(\num k)$ eventually appear within the tested bounds, forcing that
instance to have value $0$. By induction every instance is false, so the
existential is false. Substitution of a numeral can lengthen terms but
does not add logical connectives or quantifiers, so these applications of
the induction hypothesis are legitimate.

Suppose the branch were $A$-computable. Its third coordinate $e\mapsto d_e$
would then be a total $A$-computable function. Choose $q$ with
$\Phi_q^A=d$. The halting computation $\Phi_q^A(q)$ has a finite running
time and asks only finitely many oracle questions. At a sufficiently long
prefix all these questions, the position $q$ and the required running time
are available. Since $b=A$, the simulation in the membership test follows
this very computation and requires
\[
 d_q\ne\Phi_q^A(q)=d_q,
\]
a contradiction.
\end{proof}

Each of the three coordinates has now been used. The truth bits identify
the oracle against which diagonalization is to work. The witnesses make
the existential truth conditions enforceable on finite records. The
diagonal values forbid an $A$-computable completion. Without the last
coordinate the true truth table together with its least witnesses would
itself be $A$-computable.

There is a useful bound on the complexity of the particular branch in
\eqref{eq:branch-diagonal}. Let $A'$ be the halting problem relative to
$A$. This oracle can determine whether $\Phi_e^A(e)$ halts and, in the
positive case, compute its value. It also computes $A$ and the least
existential witnesses. Therefore the displayed branch is computable from
$A'$. We have given a concrete description of an infinite branch using a
stronger oracle, while proving that no branch is computable from $A$.
This does not classify the degrees of all branches.

\section{The order on natural numbers}\label{sec:order}

The subordination relation is now simple: a longer node extending a shorter
one lies below it. To make the domain literally $\N$, fix Cantor pairing
\[
 \langle a,b\rangle=\frac{(a+b)(a+b+1)}{2}+b.
\]
Encode a triple by $\langle b,\langle w,d\rangle\rangle$, and encode
finite lists recursively by
\begin{equation}\label{eq:list-code}
 \seqcode{[\,]}=0,\qquad
 \seqcode{h::s}=1+\langle h,\seqcode{s}\rangle.
\end{equation}
Decoding a nonzero list code produces a tail with a strictly smaller
code, so repeated decoding terminates. These formulas give a computable
bijection between finite lists of natural numbers and $\N$; pairing then
decodes each list entry as a triple.

Let $s_i$ be the finite record decoded from $i$. Define
\begin{equation}\label{eq:order}
 j\prec_*i
 \quad\Longleftrightarrow\quad
 s_i,s_j\in P\ \text{ and }\ s_j\text{ properly extends }s_i.
\end{equation}
If a decoded record fails the membership test, its natural number is an
isolated vertex. This permits use of every natural number as an operator
name without altering the tree's extension relation.

The direction in \eqref{eq:order} is important. The empty record, whose
code is $0$, lies above every nonempty accepted node. As records acquire
more coordinates along a branch, their codes form a descending path in
the subordination order. Their numerical magnitudes play no role in that
order.

\begin{proposition}\label{prop:order}
The relation $\prec_*$ is a recursive strict partial order. It is
ill-founded and has no $A$-computable infinite descending path.
\end{proposition}

\begin{proof}
Both node membership and proper extension of finite records are decidable.
Proper extension is irreflexive and transitive, so \eqref{eq:order} is a
recursive strict partial order.

Take the branch from Lemma~\ref{lem:branches}. The codes of its successive
prefixes, starting with the empty prefix, give an infinite descending
path. Thus the order is ill-founded.

Conversely, suppose $i_0\succ_*i_1\succ_*\cdots$ were an $A$-computable
descending path. The corresponding nodes properly extend one another,
and their lengths strictly increase. To compute the $m$th record of
their union, search along this sequence until a node of length greater
than $m$ appears and read its $m$th record. This search always terminates
and is computable from $A$. The union is consequently an $A$-computable
branch of $P$, contradicting Lemma~\ref{lem:branches}. If the sequence
starts at a nonempty node, that node simply supplies its initial records.
\end{proof}

This proposition gives the required order by an algorithm for each pair
$(j,i)$. It does not select an order from the class of all r.e. orders by
testing an uncomputable property. Nor does enumerating all of its edges
select an infinite path: such a selection would contradict the proposition.

\section{Assembling the hierarchy}\label{sec:assembly}

\begin{proof}[Proof of Theorem~\ref{thm:main}]
Use the relation \eqref{eq:order}. Proposition~\ref{prop:order} establishes
its stated order and path properties. For any index $b$ of a
closed-r.e.-generic background $B$, use the compiler of
Section~\ref{subsec:fixed-point} to produce $n$ and the family $T(n)$.
Proposition~\ref{prop:criterion} shows that this family is true. Its
generators include the particular-code schemas for every $j$ and the
truth schemas for all $j\prec_*i$, as required.

For the concrete arithmetic example, take $B_i=\PA^L$ for every $i$,
where $\PA^L$ includes induction for all formulas of the full modal
language. This family is closed-r.e.-generic by the cited arithmetic
lemmas. It is uniformly r.e., so the same compiler applies and every
$T_i$ contains this arithmetic background.

The input $b$ enters only the effective axiom and consequence enumerations
and the ensuing index compiler. The tree $P$ and order $\prec_*$ remain
fixed. Thus the successful index is obtained uniformly from $b$ for every
generic input background.
\end{proof}

The effective ingredients of the witness can be read in construction order:
\begin{enumerate}
\item Fix the arithmetic syntax list and oracle-program numbering, and
apply the finite test of Section~\ref{subsec:membership} to define $P$.
\item Decode natural numbers by \eqref{eq:list-code} and decide $j\prec_*i$
using \eqref{eq:order}.
\item Enumerate the generators for a chosen background and a proposed
index $n$, and take the finite operator closure \eqref{eq:finite-closure}.
\item Compile their consequences into the index map $f_b$ and compute
the fixed point by \eqref{eq:explicit-compiler}.
\end{enumerate}
All four are ordinary computable operations. The arithmetic truth oracle
appears in proving that a branch exists and in showing that an unsound
family would give a forbidden path. It is absent from these four algorithms.

In particular, the result specifies the entire uniformly r.e. hierarchy
at once. It does not begin by choosing a noncomputable descending path and
then try to enumerate the theories located along that path. The theories
at all natural-number labels are supplied by the common construction.

\subsection{What the members know}

Every $T_i$ contains its own particular-code schema, as well as that of
every other member. This follows because \eqref{eq:code-schema} has no
restriction relating $i$ and $j$. Truth assertions have a different
scope: \eqref{eq:truth-schema} is included only when $j\prec_*i$. In
particular, irreflexivity supplies no own-truth schema at $i$.

All code assertions refer to the consequences in the same canonical
structure in which all subordinate truth assertions are verified. No
change of interpretation along a path is used. Standard soundness of
every member follows from truth of the constructed family and the
semantic consequence relation.

This distinction between knowing a code and knowing truth is also why
the common index is harmless. It enables enumeration of what the
members conclude. It does not decide all arithmetic questions or tell a
member which downward choices can continue indefinitely.

\section{Successor chains and introspection}\label{sec:comparison}

\subsection{Why this does not give the chain in Problem 3.2.1}

Alexander's dissertation Problem~3.2.1 asks for a sequence in which
$K_i$ knows the truth and a particular code of $K_{i+1}$, together with
the Peano axioms for the full language. The companion
paper~\cite{successor} gives a negative answer in the recursively named
language and standard canonical semantics. Its supplement excludes any
computable infinite path of the corresponding endorsements.

An ill-founded relation supplies a descending sequence in the ordinary
set-theoretic sense. It need not supply an algorithm naming its
successive members. Our order demonstrates that distinction directly:
its descending paths are all noncomputable, and none can even be
computed from arithmetic truth. Relabelling one such path
$K_0,K_1,K_2,\ldots$ is therefore not a computable renaming of the
language. Uniform recursive enumerability of the full labelled family
does not make the sequence obtained by a noncomputable selection
uniformly r.e.

The proof of Theorem~\ref{thm:main} is independent of that negative
result. The absence of computable paths is proved inside the tree
construction, rather than inferred from the successor-chain obstruction.
The two papers describe compatible sides of the distinction.

\subsection{Including the minimal introspection bundle}\label{subsec:introspection}

The minimality question addressed in paper~07~\cite{minimal} concerns
background reasoning principles for the same canonical construction.
For a fixed operator $K_i$, write
\[
 \begin{aligned}
 \operatorname{AV}&=\{\phi^\eta:\phi\text{ is valid}\},\\
 i\text{-Deduction}&=
 \{\ucl(K_i(\phi\to\psi)\to(K_i\phi\to K_i\psi))\},\\
 i\text{-Introspection}&=
 \{\ucl(K_i\phi\to K_iK_i\phi)\}.
 \end{aligned}
\]
The notation $[S]_i$ places schema $S$ in the $i$th background theory and
places nothing in the other theories. The positive result from paper~07
says that, for all $i,j$,
\begin{equation}\label{eq:minimal-bundle}
 [\operatorname{AV}]_i\ \cup\
 [i\text{-Deduction}]_i\ \cup\
 [i\text{-Introspection}]_j
\end{equation}
is closed-r.e.-generic. A separate $i$-Validity supporting schema is
unnecessary for this conclusion.

The reason for the positive result is short. Let $U$ be a closed r.e.
extension of \eqref{eq:minimal-bundle}. If $U_i\vdash\rho$, compactness
gives a valid sentence
$\sigma_1\to\cdots\to\sigma_k\to\rho$ with $\sigma_\ell\in U_i$.
Assigned validity puts this valid sentence into $U_i$. Closure then
puts its $K_i$-prefix into $U_i$, and also puts each $K_i\sigma_\ell$
there. Repeated $i$-Deduction gives $U_i\vdash K_i\rho$. By the
canonical interpretation and substitution, every $i$-Introspection
instance is therefore true in $M_U$. Assigned validity and deduction
themselves are true in the canonical structure by the corresponding
source genericity facts~\cite[Lemmas~17--18]{alexander}.

The uniform union of \eqref{eq:minimal-bundle} over all pairs $(i,j)$ is
r.e. Unions of generic families are generic when the union is r.e.
\cite[Lemma~15]{alexander}. Hence that union may be added to the
arithmetic background in Theorem~\ref{thm:main}. The same tree, order
and index compiler work with the enlarged input background.

This identifies the precise role of paper~07: it permits a smaller
supporting bundle for introspection within this explicit hierarchy.
The basic witnessing hierarchy already works with the published
arithmetic background alone. The minimality assertion in paper~07 is
relative to its three named candidate supporting schemas; it does not
establish a minimal background for Theorem~\ref{thm:main}.

\section{Scope of the explicit construction}\label{sec:scope}

The order-membership test is a finite decision procedure. The index compiler
is effective and produces the r.e. theory family without consulting truth.
Existence of a descending path and soundness of the family are mathematical
theorems about those algorithms. Neither theorem requires an executable
arithmetic truth oracle to be available.

The construction retains the full quantifiers of Alexander's
Theorem~79(1): a single order works for every closed-r.e.-generic
background. Its particular-code assertions concern all members and all
formulas with free variables contained in one designated variable, as in
the source. Its subordinate truth assertions concern unrestricted modal
formulas. Full modal-language arithmetic can be included. The proof
does not restrict code assertions to arithmetic formulas or to a fixed
finite modal depth.

The language is nevertheless arithmetic with the given operators. If a
nonarithmetical predicate for computable ordinals is added, the translation
in Lemma~\ref{lem:translation} is no longer a translation into ordinary
arithmetic in the same way. The argument does not extend unchanged to
that setting. This is consistent with the separate ordinal-knowledge
well-foundedness result in Alexander's paper.

The data of the witness are a recursive order and a uniformly r.e. family.
They include no computable infinite descending sequence. A complete software
implementation of Alexander's modal consequence enumerator, and a formal
proof-assistant verification of the construction, are separate tasks. The
proof here specifies the required operations and establishes their
properties using the published logical framework.

\begin{thebibliography}{9}

\bibitem{alexander}
Samuel Allen Alexander.
\emph{Self-referential theories}.
Journal of Symbolic Logic \textbf{85}(4) (2020), 1687--1716.
\href{https://doi.org/10.1017/jsl.2020.54}{doi:10.1017/jsl.2020.54}.
Author version:
\href{https://arxiv.org/abs/2008.11535v1}{arXiv:2008.11535v1}.
Definitions 2, 7 and 12--14; Theorem 3; Lemmas 8, 10, 15, 17--18 and 20;
Theorem 24; and the arithmetic translation in the proof of Theorem 79.

\bibitem{minimal}
\emph{Minimal building blocks for introspection in self-referential theories}.
Research entry 07 in this repository.
\href{../07-minimal-introspection-bundle/minimal-introspection-bundle.pdf}
{Companion paper}.
Only its positive genericity result is used for the optional enrichment in
Section~\ref{subsec:introspection}.

\bibitem{successor}
\emph{Successor truthfulness and particular-code knowledge}.
Research entry 18 in this repository, with
\emph{Further consequences of particular-code knowledge}, Proposition S2.
\href{../18-dissertation-problem-3-2-1/successor-knowers-3-2-1.pdf}
{Companion paper}.
Used for comparison in Section~\ref{sec:comparison}, not in the construction
or proof of Theorem~\ref{thm:main}.

\end{thebibliography}
\end{document}
